\documentclass{beamer}
% imprimir
% \documentclass[handout]{beamer} 
% \usepackage{pgfpages}
% \pgfpagesuselayout{4 on 1}[a4paper,landscape,border shrink=5mm]

\mode<presentation> {
  \usetheme{Warsaw}
  \setbeamercovered{invisible}
}

\usebackgroundtemplate{\includegraphics[width=\paperwidth]{figs/libresoft-bg.png}}
\usepackage[spanish]{babel}
\usepackage[utf8]{inputenc}
\usepackage{graphics}
\usepackage{amssymb} % Simbolos matematicos

%% Metadatos del PDF.
\hypersetup{  
  pdftitle={Homogeneización del CPD},
  pdfauthor={Jose Castro, Miguel Vidal},
  pdfcreator={GSyC/Libresoft},
  pdfproducer=PDFLaTeX,
  pdfsubject={Arquitectura de servidores con software libre},
}

% Inicio de la presentación
\begin{document}

% Metadatos de la presentacion
\title{Homogeneización del CPD}
\subtitle{Arquitectura de servidores con software libre}
\institute{\{jfcastro,mvidal\}@libresoft.es} 
\author{Jose Castro, Miguel Vidal}
\date{8 de abril de 2011}

% Cover slide
\frame{
  \maketitle
  \begin{center}
    \includegraphics[width=6cm]{figs/gsyc-urjc}
  \end{center}
}

%% License slide
\begin{frame}
  \vspace{2cm}
  \begin{flushright}
    {\footnotesize \copyright{2011} Jose Castro, Miguel Vidal.} \\
    \vspace{0.25cm}
    {\scriptsize Algunos derechos reservados. \\
    Trabajo distribuido bajo una licencia Creative Commons Reconocimiento 3.0 España}\\
    \vspace{0.10cm}
  \end{flushright}
  \begin{center}
    \href{http://creativecommons.org/licenses/by/3.0/es}{\includegraphics[width=2cm]{figs/cc-by.png}} \\
    {\tiny \url{http://creativecommons.org/licenses/by/3.0/es}}
  \end{center}
\end{frame}


\section{Introducción}
%%%%%%%%%%%%%%%%%%%%%%
\begin{frame}
  \frametitle{}
  \begin{center}
    \Huge Introducción
  \end{center}
\end{frame}

\subsection{Introducción}
\begin{frame}
  \frametitle{Introducción}
  A medida que la empresa crece, va creciendo su infraestructura tecnológica:
  \begin{itemize}
    \item necesidad de almacenamiento
    \item aumenta el número de servidores
    \item diferentes sistemas operativos
    \item necesidad de nuevos servicios
    \item algunos servicios se convierten en servicios críticos
    \item ...
  \end{itemize}
\end{frame}

\begin{frame}
  \frametitle{Introducción}
  
  \begin{center}
    \alert{El resultado suele ser un conjunto heterogéneo}
  \end{center}
  \pause
  Con los problemas que esto conlleva:
  \begin{itemize}
    \item aumenta el riesgo de errores
    \item surgen problemas de escalabilidad
    \item se incrementan los agujeros de seguridad
  \end{itemize}
  \pause
  \begin{center}
    \textbf{¡La administración se convierte en una tarea compleja!}
  \end{center}
\end{frame}

\subsection{Beneficios}
\begin{frame}
  \frametitle{Beneficios}
  Una buena homogeneización del Centro de Proceso de Datos (CPD) nos va a permitir:
  \begin{itemize}
    \item reducir el tiempo de administración y despliegue
    \item aprovechar mejor los recursos hardware
    \item minimizar considerablemente el riesgo de fallo
    \item invertir en infraestructura más robusta y resistente a fallos
    \item facilitar la escalabilidad a medida de las necesidades
  \end{itemize} 
\end{frame}

\subsection{Etapas}
\begin{frame}
  \frametitle{Etapas}
  La homogeneización del CPD la podemos encarar siguiendo estas etapas:
  \begin{enumerate}
    \item Automatización de tareas
    \pause
    \item Almacenamiento como servicio
    \pause
    \item Virtualización de servidores
    \pause
    \item Alta disponibilidad de servidores y servicios
    \pause
    \item Cloud Computing -- IaaS
  \end{enumerate}
\end{frame}


\section{Automatización}
%%%%%%%%%%%%%%%%%%%%%%%%
\begin{frame}
  \frametitle{}
  \begin{center}
    \Huge Automatización
  \end{center}
\end{frame}

\subsection{Introducción}
\begin{frame}
  \frametitle{Tareas de un administrador}
  Las tareas de un administrador de sistemas se pueden clasificar en:
  \begin{itemize}
    \item Tareas repetitivas
    \item Atención a usuarios
    \item Eventos de alarmas o emergencias
    \item Despliegues y tareas nuevas
  \end{itemize}
\end{frame}

\begin{frame}
  \frametitle{Automatización}
  \begin{itemize}
    \item Sólo hay que hacer las cosas una vez
    \item Se reduce el tiempo en tareas aburridas
    \item Minimizamos el riesgo de errores
    \item Anticipación a los fallos y caídas
    \item Muchos errores se pueden solucionar automáticamente sin necesidad de intervención
  \end{itemize}
\end{frame}

\subsection{Métodos}
\begin{frame}
  \frametitle{Métodos -- Configuraciones}
  Hay que parametrizar las configuraciones en función de:
  \begin{itemize}
    \item sistema operativo
    \item distribuciones
    \item shell
    \item jerarquía de ficheros
    \item comandos
    \item servicios
    \item versiones
  \end{itemize}
\end{frame}

\begin{frame}
  \frametitle{Métodos -- Actualizaciones}
  Aviso automático de actualizaciones:
  \begin{itemize}
    \item Debian: \texttt{apticron}
    \item Red Hat / CentOS: \texttt{/etc/yum/yum-updatesd.conf} \\
      \texttt{emit\_via = email \\
      email\_to = mail@example.com \\
      email\_from = mail@example.com \\}
    \item Ports: \texttt{pkg\_version -vIL=}
  \end{itemize}
\end{frame}

\begin{frame}
  \frametitle{Métodos -- Logs}
  Herramientas para facilitar la compresión de los ficheros de log:
  \begin{itemize}
    \item Replicación de logs: \texttt{syslog-ng}
    \item Parseador de logs: scripts 
    \item Analizadores de logs: \texttt{logwatch} y \texttt{splunk}
  \end{itemize}
\end{frame}

\begin{frame}
  \frametitle{Métodos -- Sincronización}
  \textbf{NTP}
  \begin{itemize}
    \item Muchos ficheros de log en diferentes máquinas
    \item Permite tener los relojes de los servidores sincronizados
  \end{itemize}
  \vspace{0.5cm}
  \textbf{Monitorización}
  \begin{itemize}
    \item Software especializado
    \item scripts ad-hoc
  \end{itemize}
\end{frame}

\begin{frame}
  \frametitle{Métodos -- Lista daemons}
  \begin{itemize}
    \item Aviso de los sistemas y servicios por correo
    \item Lista de poco tráfico
    \item Tan sólo para mensajes de máquinas
    \item Sólo correos de errores o anomalías
    \item Buen uso de filtros: \texttt{procmail}
  \end{itemize}
\end{frame}

\begin{frame}
  \frametitle{Métodos -- Registro}
  \begin{itemize}
    \item Logging de operaciones
    \item Herramientas: \texttt{snoopy} y \texttt{rootsh}
  \end{itemize}
  \begin{itemize}
    \item Seguimiento de incidencias y operaciones
    \item Herramientas: \texttt{Trac} y \texttt{Redmine}
  \end{itemize}
\end{frame}

\begin{frame}
  \frametitle{Métodos -- Documentación}
  \begin{itemize}
    \item Política de configuración de sistemas
    \item Políticas de seguridad
    \item Políticas de actualización
    \item Procedimiento de tareas
  \end{itemize}
\end{frame}


\section{Prácticas}
%%%%%%%%%%%%%%%%%%
\begin{frame}
  \frametitle{}
  \begin{center}
    \Huge Prácticas
  \end{center}
\end{frame}

\subsection{unison}
\begin{frame}
  \frametitle{unison}
  \begin{center}
    \url{http://www.cis.upenn.edu/~bcpierce/unison/}
  \end{center}
  \begin{enumerate}
    \item Instalación de la herramienta
    \item Creación del repositorio principal
    \item Detección de ficheros a sincronizar
    \item Sincronizar desde el repositorio principal
    \item Propagación al servidor
    \item Propagación de los cambios a las máquinas
  \end{enumerate}
\end{frame}

\subsection{Scripting}
\begin{frame}
  \frametitle{Scripting}
  \textbf{Lecturas}
  \begin{itemize}
    \item ULSAH - Capítulo 2: Scripting and the shell
  \end{itemize}
  \vspace{1cm}
  \textbf{Ejercicios}
  \begin{itemize}
    \item ULSAH - Ejercicio E2.7
  \end{itemize}
\end{frame}

\subsection{Tareas repetitivas}
\begin{frame}
  \frametitle{Tareas repetitivas}
  \textbf{Lecturas}
  \begin{itemize}
    \item ULSAH - Capítulo 9: Periodic processes
  \end{itemize}
  \vspace{1cm}
  \textbf{Ejercicios}
  \begin{itemize}
    \item ULSAH - Ejercicio E9.4
  \end{itemize}
\end{frame}

\subsection{Ficheros de log}
\begin{frame}
  \frametitle{Ficheros de log}
  \textbf{Lecturas}
  \begin{itemize}
    \item ULSAH - Capítulo 11: Syslog and log files
  \end{itemize}
  \vspace{1cm}
  \textbf{Ejercicios}
  \begin{itemize}
    \item ULSAH - Ejercicio E11.4
  \end{itemize}
\end{frame}

\subsection{Mantenimiento del sistema}
\begin{frame}
  \frametitle{Mantenimiento del sistema}
  \textbf{Lecturas}
  \begin{itemize}
    \item ULSAH - Chapter 12: Software installation and management
      \begin{itemize}
        \item Managing packages
        \item Software localization and configuration: Linux and UNIX
        \item Using configuration management tools
      \end{itemize}
  \end{itemize}
  \vspace{1cm}
  \textbf{Ejercicios}
  \begin{itemize}
    \item Estudio de las herramientas \texttt{cfengine} y \texttt{puppet}
  \end{itemize}
\end{frame}

% Cover slide
\frame{
  \maketitle
  \begin{center}
    \includegraphics[width=6cm]{figs/gsyc-urjc}
  \end{center}
}

% Fin del documento
\end{document}

